% !TeX root = ../main.tex
% ---------------------------------------------------------------------------
%  Digital Underground: Geo-Integrated Models and Digital Twins
%  Enriched from the author's study report on sustainability and resilience of
%  underground infrastructure, which supplies the geo-integrated asset model
%  and the four-level digital twin framework used to organise this chapter.
%  Other sources: buildingSMART IFC-Tunnel; ISO 19650; ITA WG22 on BIM.
%  Related work in the parent repository: CRC_Book_Chapter/BIM_DigitalTwin,
%  Tunnel-DT2026, TunVision, Digital-Twin.
% ---------------------------------------------------------------------------
\chapter{Digital Underground: Geo-Integrated Models and Digital Twins}
\label{ch:digital-underground}

A geometric model of a tunnel is of limited engineering value. What underground
projects need is a model that also carries the ground: strata and inferred
boundaries, hydrogeological units, design parameters and their uncertainty,
construction sequence and support timing, quantities and carbon factors,
monitoring and inspection records. That object --- the geo-integrated asset
model --- is what allows design, construction, carbon accounting, risk
management and operation to draw on the same project information rather than
maintaining separate and divergent versions of it.

This chapter treats digital practice as information flow and decision
traceability, not as software or visualisation.

\section{Digital transformation and information management}
\label{sec:digital-transformation-and-information-management}

\subsection{What has actually changed}
\label{sec:digital-transformation}
\scopenote{Demonstrated capability against routine adoption, and the
organisational rather than technical obstacles to closing that gap.}

\subsection{BIM for linear underground assets}
\label{sec:bim-underground}
\scopenote{Why building-oriented BIM adapts poorly to alignments and to
ground: parametric alignment modelling, the ground as a modelled object rather
than a backdrop, and level of information need.}

\section{Standards and data environments}
\label{sec:standards-and-data-environments}

\subsection{Data standards and interoperability}
\label{sec:data-standards}
\scopenote{IFC and the IFC-Tunnel extension, LandXML, ISO 19650 information
management, and the cost of proprietary lock-in over an asset life measured in
a century.}

\subsection{Common data environments}
\label{sec:cde}
\scopenote{Single source of truth, version and status control, and the
contractual standing of model information.}

\section{The geo-integrated asset model}
\label{sec:geo-integrated-model}

\subsection{What the model must contain}
\label{sec:model-contents}
The model is useful only when it integrates, in one structure:

\begin{itemize}[nosep]
  \item infrastructure function and system interfaces;
  \item alignment, geometry, lining, support, shafts, stations and fit-out;
  \item borehole data, geological interpretation, hydrogeological units,
        design parameters and uncertainty;
  \item construction method, sequence, support timing, machine performance and
        haulage assumptions;
  \item quantities, carbon factors, energy factors and the declared system
        boundary;
  \item monitoring, inspection, condition, damage and maintenance records.
\end{itemize}

\subsection{Ground models as digital objects}
\label{sec:digital-ground-model}
\scopenote{Three-dimensional geological modelling, implicit surface methods,
chainage-based ground classification, and the problem of representing
uncertainty inside a model that renders as though it were certain.}

\subsection{Level of detail as a decision variable}
\label{sec:level-of-detail}
\scopenote{Detail should follow the decision being made. Early planning needs
system-level quantities; detailed design, construction control and operation
need component- and material-level information. Modelling everything to the
finest level is a common and expensive error.}

\section{The analysis loop}
\label{sec:analysis-loop}

\subsection{Instantiating analysis models}
\label{sec:model-instantiation}
\scopenote{Generating settlement, support-interaction, lining and
excavation-sequence models from the asset and ground models rather than
rebuilding geometry by hand for each analysis.
Chapter~\ref{ch:advanced-numerical}.}

\subsection{Forward prediction and backward updating}
\label{sec:forward-backward}
\scopenote{Forward analysis predicts performance; observation and validation
data feed back to update parameters, refine the design or change construction
control. The loop must stay active through construction and operation, which
is what distinguishes a working model from a drawing container.}

\section{Digital twins}
\label{sec:digital-twins}

\subsection{Four levels of twin maturity}
\label{sec:twin-maturity}
The useful distinction between a model and a twin is the feedback loop. Four
levels can be separated:

\begin{description}[leftmargin=28mm,style=nextline,nosep]
  \item[Descriptive] What is the asset? Geometry, geology, systems, materials
    and records organised into a structured asset model.
  \item[Reflective] What is happening now? Monitoring, machine data, inspection
    and condition observations validate the model and diagnose current state.
  \item[Predictive] What will happen next? Physical, statistical, machine
    learning and surrogate models forecast settlement, convergence, leakage,
    deterioration, damage, energy, carbon and functionality.
  \item[Prescriptive] What should be done? Intervention options are compared on
    risk, recovery, cost, time, carbon and service outcome.
\end{description}

\subsection{Closing the loop}
\label{sec:twin-loop}
\scopenote{Each observation and intervention updates the model for the next
decision. Without that return path, a twin is a dashboard.}

\section{Analytics, prediction and governance}
\label{sec:analytics-and-governance}

\subsection{Construction and machine data analytics}
\label{sec:tbm-analytics}
\scopenote{Machine data streams, cleaning and ring-based aggregation, and
inference of ground conditions from machine response ahead of and behind the
face.}

\subsection{Machine learning and automated inspection}
\label{sec:machine-learning}
\scopenote{Supervised prediction of penetration rate, ground class and
settlement; automated photogrammetric and LiDAR face mapping; lining defect
detection. The transferability problem between projects, and the danger of
models that extrapolate confidently outside their training envelope.}

\subsection{Data quality, governance and responsibility}
\label{sec:data-governance}
\scopenote{Model provenance, auditability, and where engineering judgement and
accountability sit when a model informs a safety-critical decision. Emerging
methods should be presented as such until supported by mature project
evidence.}
